% Part: normal-modal-logic
% Chapter: axioms-systems
% Section: soundness

\documentclass[../../../include/open-logic-section]{subfiles}

\begin{document}

\olfileid[pt]{nml}{prf}{snd}

\olsection{Correção}

!!^a{derivation} é dito correto se tudo o que pode ser !!{derive}d é
válido. Ao considerar sistemas modais, O.s.\@, !!{derivation}s nas
quais, além de \Ax{K}, podemos usar instâncias de certas !!{formula}s
$!A_1$, \dots,~$!A_n$, queremos que toda fórmula !!{derivable} seja
verdadeira em qualquer modelo no qual $!A_1$, \dots, $!A_n$ sejam
verdadeiras.

\begin{thm}[Teorema da correção]\ollabel{thm:soundness}
  Se toda instância de $!A_1$, \dots, $!A_n$ é válida nas classes de
  modelos $\mClass{C}_1$, \dots, $\mClass{C}_n$, respectivamente,
  então $\Log{K}!A_1\dots !A_n \Proves !B$ implica que $!B$ é válida
  na classe de modelos
  $\mClass{C}_1 \cap \dots \cap \mClass{C}_n$.
\end{thm}

\begin{proof}
  Por indução no comprimento das demonstrações. Para abreviar, ponha $\mClass{C} =
  \mClass{C}_1 \cap \dots \cap \mClass{C}_n$.
  \begin{enumerate}
  \item Base da indução: Se $!B$ tem uma demonstração de
    comprimento~$1$, então ela é ou uma instância tautológica, ou uma
    instância de~\Ax{K},\iftag{prvDiamond}{ ou de~\Dual{},}{} ou uma
    instância de um dos axiomas $!A_1$, \dots,~$!A_n$. No primeiro
    caso, $!B$ é válida em $\mClass{C}$, pois instâncias tautológicas
    são válidas em \emph{qualquer} classe de modelos, por
    \olref[syn][tau]{prop:valid-taut}. O mesmo vale no segundo caso,
    por \olref[syn][sch]{prop:Kvalid}\iftag{prvDiamond}{ e
      \olref[syn][sch]{prop:Dual-valid}}{}. Finalmente, no terceiro
    caso, como $!B$ é válida em $\mClass{C}_i$ e $\mClass{C} \subseteq
    \mClass{C}_i$, temos que $!B$ também é válida em $\mClass{C}$ por
    \olref[syn][val]{prop:subset-class}.
  \item Passo indutivo: Suponha que $!B$ tenha uma demonstração de
    comprimento $k>1$. Se $!B$ é uma instância tautológica ou uma
    instância de um dos axiomas $!A_1$, \dots, $!A_n$, procedemos como
    no passo anterior. Suponha, portanto, que $!B$ seja obtida por
    \MP{} a partir das !!{formula}s anteriores $!C \lif !B$ e $!C$.
    Então $!C \lif !B$ e $!C$ têm demonstrações de comprimento $<k$ e,
    pela hipótese de indução, são válidas em~$\mClass{C}$. Por
    \olref[syn][sch]{prop:soundMP}, $!B$ também é válida em
    $\mClass{C}$. Finalmente, suponha que $!B$ seja obtida por \Nec{}
    a partir de $!C$ (de modo que $!B = \Box!C$). Pela hipótese de
    indução, $!C$ é válida em $\mClass{C}$, e por
    \olref[syn][val]{prop:Nec-rule}, $!B$ também o é.
  \end{enumerate}
\end{proof}

\end{document}
